\documentclass[11pt]{article}
\usepackage{preamble}
\setlength{\parskip}{4pt}

\begin{document}

\daybanner{AP Statistics \textperiodcentered\ Unit 1 \textperiodcentered\ Topic 1.10 \textperiodcentered\ B: 2026-09-25 \textperiodcentered\ E: 2026-10-02 \textperiodcentered\ TEACHER KEY}{Who Chose Whom?}

\sectionbanner{CED TETHER}

{\small
\begin{itemize}[leftmargin=1.5em,itemsep=2pt,topsep=2pt]
  \item \textbf{VAR-1.E} Identify questions to be answered about data collection methods. \textbf{[Skill 1.A]}
  \item \textbf{VAR-1.E.1} Methods for data collection that do not rely on chance result in untrustworthy conclusions.
\end{itemize}
}
\textbf{Essential question:} How does the way people enter a sample limit which result we should trust?\par
\textbf{DOK-3 skill:} 4.A \textperiodcentered\ \textbf{FRQ pattern:} {\small\texttt{voluntary-\allowbreak{}response-\allowbreak{}vs-\allowbreak{}random-\allowbreak{}which-\allowbreak{}number-\allowbreak{}and-\allowbreak{}why}}\par
\textbf{DOK rationale:} Part (c) coordinates selection method, response count, and target population to choose between competing estimates, explain why sample size cannot repair self-selection, and match a design change to the bias.\par

\frameworkphaseheader{Do Now (first take) $\rightarrow$ Explore (video + follow-along) $\rightarrow$ Exit (finish + turn in)}{1 $\rightarrow$ 3}{45}{%
\begin{itemize}[leftmargin=1.5em,itemsep=2pt,topsep=2pt]
  \item Board slide up at the bell. Read both headlines without endorsing either.
  \item Begin the video follow-along at minute 5; leave first takes ungraded.
  \item Return to the board slide at minute 33. Circulate and ask who chose whom in each poll.
  \item Collect at the door. Score part (c) only, E/P/I.
\end{itemize}
}{%
\begin{itemize}[leftmargin=1.5em,itemsep=2pt,topsep=2pt]
  \item Commit to one headline and one supporting detail before the lesson.
  \item Complete the video follow-along about trustworthy data collection.
  \item Finish (a)--(c), revise the first take in the exit line, and turn in the sheet.
\end{itemize}
}{%
\begin{itemize}[leftmargin=1.5em,itemsep=2pt,topsep=2pt]
  \item In which method did students choose the sample, and in which did the researcher choose it?
  \item What kind of student is especially likely to click the website poll?
  \item What does a larger response count improve, and what can it not repair?
  \item How would your fix stop students from volunteering themselves into the sample?
\end{itemize}
}{Solo teacher. Circulate ~5 pairs. Require students to point to how each sample was formed before accepting a percentage; do not say which headline is stronger first.}

\begin{callout}{\IconWarn\ WATCH FOR}{calloutred}
\begin{itemize}[leftmargin=1.5em,itemsep=2pt,topsep=2pt]
  \item Treating the larger response count as proof that the website sample represents the school.
  \item Calling the homeroom method an SRS of students; the random units are homerooms.
  \item Proposing a wording change that leaves the website poll open to self-selection.
\end{itemize}
\end{callout}

\sectionbanner{SCORING PART (c) --- E / P / I}

\textbf{Expected elements}
\begin{itemize}[leftmargin=1.5em,itemsep=2pt,topsep=2pt]
  \item \textbf{chooses-random-homeroom-estimate} (required) --- Chooses 44 percent and ties the choice to chance selection of homerooms plus asking every student in them
  \item \textbf{diagnoses-self-selection} (required) --- Explains that website visitors choose whether to respond, so strong opinions may be overrepresented
  \item \textbf{rejects-large-n-repair} (required) --- Explains that 412 responses do not make a self-selected sample representative
  \item \textbf{proposes-chance-based-fix} (required) --- Proposes a chance-based selection from the school population and prevents open self-selection
\end{itemize}

{\small\renewcommand{\arraystretch}{1.25}
\noindent\begin{tabular}{|p{0.5in}|p{5.9in}|}
\hline
\textbf{E} & Chooses the homeroom estimate, cites its chance-based selection, explains both the website's self-selection and why 412 does not repair it, and gives a workable chance-based fix. \\ \hline
\textbf{P} & Makes the supported choice and identifies self-selection but omits either the large-sample limitation or a workable fix, or gives a fix without explaining how it changes selection. \\ \hline
\textbf{I} & Chooses 71 percent because 412 is larger, treats both polls as equally representative, or gives no selection-method evidence. \\ \hline
\end{tabular}}

\textbf{Common mistakes}
\begin{itemize}[leftmargin=1.5em,itemsep=2pt,topsep=2pt]
  \item Choosing the larger sample solely because 412 exceeds 68
  \item Calling the website poll random because many different visitors could answer
  \item Saying only that the polls disagree without analyzing who chose to respond
  \item Keeping the poll open to everyone but changing only its wording
\end{itemize}

\clearpage
\focusbanner{TODAY'S PROBLEM --- annotated}

Suppose a school website posts the question, ``Should the cafeteria go cashless?'' Any visitor may answer. Of 412 responses, $71\%$ say Yes. That same week, a teacher randomly selects three homerooms and asks every student in those rooms. Of 68 responses, $44\%$ say Yes.\par\smallskip The website headline is ``Students back a cashless cafeteria.'' The teacher's headline is ``Students are split on going cashless.''\par
\begin{center}
\textbf{Two polls of student opinion}\par\smallskip
\begin{tabular}{|l|c|c|c|}
\hline
\textbf{Poll} & \textbf{$n$} & \textbf{Yes} & \textbf{Selection} \\ \hline
\textbf{Website} & 412 & $71\%$ & Visitors opt in \\ \hline
\textbf{Homerooms} & 68 & $44\%$ & 3 rooms random \\ \hline
\end{tabular}
\end{center}
\begin{firsttakebox}{FIRST TAKE (5 min)}
Which headline would you trust more for what students at this school think? Name one detail behind your choice.\par
\answer{Not graded. Expect many students to choose 71 percent because 412 is larger.}
\end{firsttakebox}

\textit{Rules callout ``TRUST STARTS WITH SELECTION (keep on the page)'' is printed on the student sheet.}\par\medskip

\dokbadge{1}~\textbf{(a)} State the question of interest and the population the two headlines are trying to describe.\par
\answer{Question: What proportion of students at this school support a cashless cafeteria? Population: all students at the school.}

\dokbadge{2}~\textbf{(b)} For each poll, describe how people entered the sample. Name the sampling problem, if any.\par
\answer{The website poll is a voluntary-response sample because site visitors choose whether to answer. The homeroom poll uses chance to select three homerooms and then asks every student in those rooms; it is a cluster sample of homerooms.}

\dokbadge{3}~\textbf{(c)} Which percentage should a reader trust more for what students at this school think? Cite the selection feature that settles your choice. Explain why 412 responses do not repair the other poll, and propose one change that would make the website poll more trustworthy.\par
\answer{The $44\%$ homeroom result is more trustworthy because chance selected the homerooms and every student in them was asked. The $71\%$ website result is vulnerable to voluntary-response bias: students with strong views or who visit the site choose themselves. Its 412 responses reduce random sampling variability only for the self-selected group; they do not make that group representative. A fix is to randomly select students from the full roster and invite only those selected, following up with nonresponders. Randomly selecting homerooms and asking everyone, as the teacher did, is also acceptable.}

\begin{summaryexitbox}{EXIT}
One thing the video changed about my first take: \hrulefill
\end{summaryexitbox}

\end{document}
